% ===================================================================
%  LENTELĖ Nr. 11 — Miestas ir jo paminklai. --- Gaisras.
%  Traduction 2026 d'apres texto/io, controlee sur texto/fr, texto/ru
%  et texto/cs. Consigne : voir texto/lt/10-lentele-01.tex.
%
%  LA VILLE EST, POUR LE GALICIEN, L'UN DES ENDROITS OU LA DISTANCE A
%  LA LANGUE VOISINE EST MINIMALE. POUR LE LITUANIEN C'EST L'ENDROIT
%  OU LA LANGUE A FAILLI NE PLUS ETRE. Jusqu'en 1918 Vilnius et
%  Kaunas parlaient polonais et russe ; l'imprime lituanien en
%  caracteres latins a ete interdit de 1864 a 1904, et les livres
%  passaient la frontiere sur le dos des « knygnešiai », les
%  porte-livres. Le vocabulaire administratif d'une ville lituanienne
%  a donc ete fait, pour l'essentiel, entre 1918 et 1940.
%
%  ET IL A ETE FAIT COMME LE TCHEQUE AVAIT FAIT LE SIEN, avec les
%  memes deux mots. La colonne treize avait montre le reveil national
%  tcheque decidant contre l'allemand : « spořitelna » la caisse
%  d'epargne, « hasiči » les pompiers, litteralement les eteigneurs.
%  Le lituanien a « taupomoji kasa », la caisse qui epargne, et
%  « ugniagesiai », les eteigneurs-de-feu. DEUX REVEILS, DEUX
%  VOISINS DIFFERENTS — l'allemand la, le polonais et le russe ici —
%  ET LES DEUX MEMES MOTS REFAITS DE LA MEME FACON.
%
%  MAIS LA COMPARAISON NE VA PAS JUSQU'AU BOUT, ET C'EST LA QU'ELLE
%  DEVIENT INSTRUCTIVE. Le tcheque avait refuse « rathaus » et forge
%  « radnice » ; le lituanien a garde « rotušė », qui est le Rathaus
%  allemand passe par le polonais. La meme langue qui refait la
%  caisse d'epargne et les pompiers laisse a l'hotel de ville son nom
%  etranger. Une langue qui decide ne decide pas partout, et le
%  tableau 11 montre qu'elle ne decide pas non plus selon
%  l'importance de la chose : la mairie compte plus que la caisse
%  d'epargne.
%
%  ET « ligoninė » EST L'HOPITAL, DE « liga » LA MALADIE : le suffixe
%  de lieu du tableau 6 — « virtuvė », « skalbykla » — applique cette
%  fois a un batiment public. Ce que la langue fabriquait pour une
%  maison, elle le fabrique ici pour une ville.
%
%  DEUX GENITIFS TENUS PAR UNE RELATIVE ET UNE APPOSITION : la
%  coupole (40) de l'hopital (39), les musiciens (81) de l'orchestre
%  (80).
% ===================================================================

%%K t11-apar-1 apar
\VUsaut{2.0mm}
\VUcentre{13.2pt}{90}{TREČIOJI SERIJA}
\VUsaut{3.0mm}
\VUfilet{16mm}
\VUsaut{5.6mm}
\VUtitre{15.0pt}{60}{LENTELĖ Nr. 11}
\VUsaut{1.4mm}
\VUpk{11.4pt}{40}{Miestas ir jo paminklai. --- Gaisras.}
\VUsaut{2.2mm}

%%K t11-tit-1 sub
\VUpk{10.2pt}{40}{Priemiestis.}
\VUsaut{3.0mm}

%%K t11-c1-01-1 p
1. --- Gyvenu dideliame mieste, esančiame už 100 kilometrų nuo Vandenyno ir vis dėlto
esančiame aukšto rango \VUgras{uostu} \textsuperscript{(1)}. Upė, kuri jį juosia, yra
400 metrų pločio ir sudaro labai gražų reidą.

%%K t11-c1-02-1 p
2. --- Visa miesto dalis, esanti prie \VUgras{krantinių} \textsuperscript{(2)},
atrodo visai jūriška ir pramoninė ir sudaro \VUgras{priemiestį}
\textsuperscript{(3)}, kuriame gyvena tik laivininkai ir darbininkai. Šie dirba
\VUgras{fabrikuose} \textsuperscript{(4)} ir \VUgras{dujų fabrike} \textsuperscript{(5)},
kurio aukštas \VUgras{kaminas} \textsuperscript{(6)} ir \VUgras{dujų rezervuaras}
matyti iš labai toli.

%%K t11-c1-03-1 p
3. --- Viduramžiais miestas buvo įtvirtintas, ir dar randama keletas jo pylimų
liekanų, ypač \VUgras{vartai} \textsuperscript{(8)}, likę iš senosios \VUgras{įtvirtintos
pilies} \textsuperscript{(7)} ir iš šonų remiami dviejų \VUgras{bokštų}
\textsuperscript{(9)},
kurie dabar naudojami kaip \VUgras{kalėjimas.}

%%K t11-c1-04-1 p
4. --- Visame tame \VUgras{kvartale,} kuris atrodo labai senas, tik vienas statinys
palyginti naujas: tai \VUgras{muitinė} \textsuperscript{(10)}. Ji stovi tarp krantinės
ir \VUgras{gatvelės} \textsuperscript{(11)}, kuri veda prie senosios \VUgras{rotušės}
\textsuperscript{(12)}. Tą paskutinį paminklą lengva atpažinti iš milžiniškos
\VUgras{varpinės} \textsuperscript{(13)}, kuri dominuoja virš senojo miesto.

%%K t11-c1-05-1 p
5. --- Kitoje gatvės pusėje \VUgras{stovi} pastatas, pastatytas tuo pačiu stiliumi
kaip muitinė: tai \VUgras{Birža} \textsuperscript{(14)}. Vienas jos fasadas žiūri į
mažą aikštę, kurioje yra \VUgras{apskrities valdyba} \textsuperscript{(15)}. Iš tikrųjų
mūsų miestas, nebūdamas sostine, vis dėlto yra svarbus apskrities centras.

%%K t11-tit-2 sub
\VUsaut{5.0mm}
\VUpk{10.2pt}{40}{Aukščiausioji miesto dalis.}
\VUsaut{3.4mm}

%%K t11-c1-06-1 p
6. --- Įgula, gana gausi, apgyvendinta plačiose ir labai sveikose
\VUgras{kareivinėse} \textsuperscript{(16)}; jos tęsiasi iki \VUgras{miesto parko}
\textsuperscript{(17)} grotų. \VUgras{Bulvaras} \textsuperscript{(19)}, kuris veda į tą
parką, eina už \VUgras{taupomosios kasos} \textsuperscript{(18)} ir baigiasi
\VUgras{katedros} \textsuperscript{(20)} aikštėje.

%%K t11-c1-07-1 p
7. --- Visi tie paminklai išsidėstę amfiteatru ant kalvelės, kurioje yra ir mūsų
plati \VUgras{licėjus} \textsuperscript{(21)}.

%%K t11-c1-07-2 p
Jeigu leidiesi kiek nuožulniu keliu, kuris išeina į Didžiąją gatvę, prieini
\VUgras{Teismą} \textsuperscript{(22)}, prieš kurį plyti malonus \VUgras{viešasis sodas}
\textsuperscript{(26)}. Tas \VUgras{skveras} nelabai didelis, bet miesto viduryje jis
sudaro žalumos ir vėsos oazę. Todėl jį labai lanko miestiečiai, kurie mėgsta ateiti
atsisėsti po \VUgras{kavinės} \textsuperscript{(25)} tentu ir klausytis kareivių
muzikantų, kurie ketvirtadieniais ir sekmadieniais groja \VUgras{paviljone}
\textsuperscript{(27)}, pastatytame tarp vejų ir gėlynų.

%%K t11-c1-08-1 p
8. --- Ant vienos iš tų vejų stovi bronzinė \VUgras{statula} \textsuperscript{(28)},
paminklas, pastatytas atminti vieną iš mūsų miestiečių, kuris padarė didelių paslaugų
savo gimtajam miestui.

%%K t11-tit-3 sub
\VUsaut{4.0mm}
\VUpk{10.2pt}{40}{Susisiekimo priemonės.}
\VUsaut{3.4mm}

%%K t11-c1-09-1 p
9. --- Iki šiol mūsų miestas dar neapšviestas elektros šviesa, jis teturi
\VUgras{dujų degiklius} \textsuperscript{(29)}. Bet \VUgras{vietos spauda} veda
kampaniją, kad ta apšvietimo sistema būtų pagerinta, \VUgras{kaip} jau buvo
patobulinta \VUgras{tramvajų traukos sistema.} \VUgras{Senieji} vagonai,
\VUgras{traukiami} \VUgras{arklių,} praktiškai pakeisti \VUgras{elektriniais
tramvajais} \textsuperscript{(31)}, kurie greitai perveža keleivius iš vieno
\VUgras{miesto galo į kitą} už labai \VUgras{pigią} kainą.

%%K t11-c1-10-1 p
10. --- Prie Teismo rūmų yra \VUgras{Bankas} \textsuperscript{(32)}, kuriame
diskontuojami prekybos vertybiniai popieriai. \VUgras{Banko pasiuntiniai}
\textsuperscript{(33)} eina į namus išieškoti skolų.

%%K t11-tit-4 sub
\VUsaut{4.0mm}
\VUpk{10.2pt}{40}{Menas ir švietimas.}
\VUsaut{3.4mm}

%%K t11-c1-11-1 p
11. --- Gražus pastatas, kuris stovi netoliese, yra \VUgras{muziejus.} Jame kartu yra
miesto \VUgras{biblioteka} \textsuperscript{(34)}, gražūs \VUgras{gamtos mokslų
rinkiniai} \textsuperscript{(35)} (Gamtos muziejus) ir grakštus \VUgras{paveikslų
muziejus} \textsuperscript{(36)}, kuris sudaro puikią nuolatinę \VUgras{parodą.}

%%K t11-c1-11-2 p
Jis atidaromas visuomenei du kartus per savaitę, bet svetimšaliai gali jį lankyti bet
kurią dieną už mažą arbatpinigį \VUgras{sargui} \textsuperscript{(37)}.
\VUgras{Dailininkai} \textsuperscript{(38)} ateina ten dirbti. Matome juos
\VUgras{prieš} savo molbertus, su \VUgras{palete} rankoje, kopijuojančius garsius
paveikslus.

%%K t11-c1-12-1 p
12. --- Ant aukštumos, už muziejaus, pastebime \VUgras{kupolą} \textsuperscript{(40)},
kuris dengia \VUgras{ligoninę} \textsuperscript{(39)}, ir \VUgras{mokytojų seminarijos}
\textsuperscript{(41)} pastatus, kurioje buvo rengiami mūsų valsčiaus mokyklų
mokytojai.

Teatro aikštėje yra \VUgras{pašto kontora} \textsuperscript{(43)}, įrengta
\VUgras{privačiame name} \textsuperscript{(42)}.

%%K t11-tit-5 sub
\VUsaut{5.0mm}
\VUpk{11.4pt}{40}{Gaisras.}
\VUsaut{3.0mm}

%%K t11-tit-6 sub
\VUpk{10.2pt}{40}{Šauksmas dėl gaisro.}
\VUsaut{3.4mm}

%%K t11-c2-01-1 p
1. --- Baisi žinia sklinda per miestą: teatre ką tik kilo gaisras! Tą akimirką, kai
atbėgu į \VUgras{aikštę} \textsuperscript{(96)}, minia jau susirinkusi ant
\VUgras{šaligatvio} \textsuperscript{(46)} ir ant \VUgras{grindinio}
\textsuperscript{(47)}. \VUgras{Paštininkai} \textsuperscript{(44)} ir
\VUgras{laiškininkai} \textsuperscript{(45)} išeina iš pašto kontoros.
\VUgras{Tramvajaus vairuotojas} \textsuperscript{(48)} turėjo sustabdyti savo vagoną,
kad praleistų miesto \VUgras{greitosios pagalbos vežimą} \textsuperscript{(50)}, kuris
atvyksta su \VUgras{chirurgu} \textsuperscript{(52)} ir \VUgras{slaugytoju}
\textsuperscript{(51)} slaugyti \VUgras{sužeistojo} \textsuperscript{(53)}, atnešto
\VUgras{neštuvais} \textsuperscript{(54)} prie \VUgras{laikraščių kiosko}
\textsuperscript{(55)}.

%%K t11-c2-02-1 p
2. --- Pėstininkų kuopa, kuri grįžo iš pratybų, sustojo padėti palaikyti tvarką kartu
su \VUgras{raitaisiais miesto žandarais} \textsuperscript{(61)}. \VUgras{Raiteliai,}
kurių arkliai stojasi piestu, geriau negu \VUgras{kareiviai} \textsuperscript{(57)} ar
\VUgras{žandarai} \textsuperscript{(63)} atstumia \VUgras{smalsuolius}
\textsuperscript{(56)}. Beje, žandarai labai užsiėmę. Vienas iš jų rodo
\VUgras{policininkams} \textsuperscript{(64)}, kur reikia nugabenti kitą
\VUgras{nukentėjusįjį} \textsuperscript{(65)}, drąsų ugniagesį, nukritusį nuo
kopėčių.

%%K t11-c2-03-1 p
3. --- \VUgras{Karininkas} \textsuperscript{(66)} tiksliai nurodo vietą, kur reikia
pastatyti atvažiuojantį \VUgras{garinį siurblį} \textsuperscript{(67)}. Laukdamas tos
galingos mašinos, jis skubiai liepė pastatyti \VUgras{rankinį siurblį}
\textsuperscript{(68)}, kurio ilga \VUgras{žarna} \textsuperscript{(69)} lieja vandenį
srautu ant ugnies.

Šeši ugniagesiai jį valdo, o jų draugas, kuris laiko \VUgras{švirkštą}
\textsuperscript{(70)}, kreipia \VUgras{čiurkšlę} \textsuperscript{(71)} į gaisro
vidurį.

%%K t11-c2-04-1 p
4. --- Kitoje teatro pusėje pastatytos didelės \VUgras{kopėčios}
\textsuperscript{(72)}. Jos leidžia ugniagesiams prisiartinti prie liepsnų, kurios
trykšta tarp juodų \VUgras{dūmų} \textsuperscript{(75)} sūkurių.

%%K t11-tit-7 sub
\VUsaut{5.0mm}
\VUpk{10.2pt}{40}{Narsūs žygiai.}
\VUsaut{3.4mm}

%%K t11-c2-05-1 p
5. --- \VUgras{Gaisras} \textsuperscript{(74)} kilo fojė, \VUgras{teatro}
\textsuperscript{(73)} prieškambaryje, kai buvo repetuojama \VUgras{opera,} kurios
pirmasis \VUgras{spektaklis} turėjo įvykti rytoj. Laimei, žiūrovų salėje tebuvo
nedaug \VUgras{žiūrovų} \textsuperscript{(76)}. Vargšai prisiglaudė ant
\VUgras{balkono} \textsuperscript{(77)}, kur drąsūs \VUgras{ugniagesiai}
\textsuperscript{(86)} ateina jų gelbėti su kopėčiomis ir lynais.

%%K t11-c2-06-1 p
6. --- Po \VUgras{peristiliu} \textsuperscript{(78)}, kur \VUgras{afiša}
\textsuperscript{(79)} skelbia apie spektaklį, matome bėgančius \VUgras{muzikantus}
\textsuperscript{(81)}, \VUgras{orkestro} \textsuperscript{(80)} narius, kurie
išsineša savo instrumentus: \VUgras{smuiką} \textsuperscript{(81)},
\VUgras{fleitą} \textsuperscript{(82)}, \VUgras{violončelę} \textsuperscript{(83)},
\VUgras{trombonu} \textsuperscript{(84)}, \VUgras{būgną} \textsuperscript{(85)} ir
t.~t. Stengiamasi išgelbėti dalį \VUgras{baldų} \textsuperscript{(87)}.

%%K t11-c2-07-1 p
7. --- \VUgras{Aktoriai} \textsuperscript{(89)} ir \VUgras{aktorės}
\textsuperscript{(88)}, \VUgras{dainininkai} ir \VUgras{dainininkės} turėjo bėgti su
savo teatro \VUgras{kostiumais} \textsuperscript{(90)}. Tenoras aiškina
\VUgras{žurnalistui} \textsuperscript{(92)}, kaip tai atsitiko. \VUgras{Reporteris}
tuoj pat atbėgo nuo pirmos akimirkos su valdžios vyrais: \VUgras{apskrities
viršininku} \textsuperscript{(91)}, \VUgras{burmistru} \textsuperscript{(93)},
\VUgras{policijos komisaru} \textsuperscript{(94)} ir \VUgras{teismo pareigūnu}
\textsuperscript{(95)}, kurie ves tyrimą, kad nuspręstų dėl atsakomybės.

\VUsaut{6.0mm}
\VUfilet{22mm}
